\documentclass[10pt,UTF8]{article}

\usepackage{geometry}

\usepackage{ctex}

\geometry{a4paper,scale=0.9}
\ctexset{today=old}

\usepackage[bookmarksnumbered=true]{hyperref}  % 生成大纲


\usepackage{times}  % 设置正文字体为 Adobe Times Roman

% \usepackage{makecell}
% \usepackage{multirow} % 表格多行合并

% \usepackage{graphicx} %插入图片的宏包
% \usepackage{subfigure}
% \usepackage{float} %设置图片浮动位置的宏包

\usepackage{amsmath}
\usepackage{amssymb}
\usepackage{mathtools}

\usepackage{physics}

% \usepackage{siunitx}
% % 关闭警告
% \ExplSyntaxOn
% \msg_redirect_name:nnn { siunitx } { physics-pkg } { none }
% \ExplSyntaxOff

%%%%%%%%%%%%%%%%%%%%%%%% tcolorbox %%%%%%%%%%%%%%%%%%%%%%%%%%
\usepackage[most]{tcolorbox}

\newenvironment{definition box}
{\begin{tcolorbox}[colback=blue!5!white,colframe=blue!75!black,parbox=false,before upper=\par,before lower=\par]}
{\end{tcolorbox}}

\newenvironment{theorem box}
{\begin{tcolorbox}[colback=yellow!5!white,colframe=yellow!75!black,parbox=false,before upper=\par,before lower=\par]}
{\end{tcolorbox}}
%%%%%%%%%%%%%%%%%%%%%%%%%%%%%%%%%%%%%%%%%%%%%%%%%%%%%%%%%%%%
            
%%%%%%%%%%%%%%%%%%%%% amsthm %%%%%%%%%%%%%%%%
\usepackage{amsthm}

\theoremstyle{definition}
\newtheorem{definition inner}{Definition}[section]
\newenvironment{definition}[1][]
{\begin{definition box}\begin{definition inner}[#1]\hfill\par}
{\end{definition inner}\end{definition box}}

\theoremstyle{plain}
\newtheorem{theorem inner}{Theorem}[section]
\newenvironment{theorem}[1][]
{\begin{theorem box}\begin{theorem inner}[#1]\hfill\par}
{\end{theorem inner}\end{theorem box}}

\theoremstyle{definition}
\newtheorem*{proof inner}{\textit{Proof}}
\renewenvironment{proof}[1][]
{\begin{proof inner}[#1]\hfill\par}
{\qed\end{proof inner}}

\theoremstyle{remark}
\newtheorem*{remark}{\textit{Remark}}
%%%%%%%%%%%%%%%%%%%%%%%%%%%%%%%%%%%%%%%%%%%%%%


\newcommand{\mycases}[2][0.8]{
    \left\{\vphantom{\scalebox{#1}{
        $\begin{matrix}#2\end{matrix}$
    }}\right.\!\!\!\!
    \begin{array}{l@{\quad}l}#2\end{array}
}

\newcommand{\ju}{\mathrm{j}}  % 虚数单位
\newcommand{\e}{\mathrm{e}{\mkern 1 mu}}


\begin{document}

\part*{Fourier Series Representation of Periodic Signals}

\section{The Response of LTI Systems to Complex Exponentials}

\begin{definition}[Eigenfunctions \& Eigenvalues]
    A signal for which the system output is a
    constant times the input is referred to as an \emph{eigenfunction} of the system,
    and the amplitude factor is referred to as the system's \emph{eigenvalue}.
\end{definition}

It is easy to show that \emph{complex exponentials} are eigenfunctions of LTI systems.

\begin{equation*}
    \e^{s t} \xlongrightarrow{\text{a CT LTI system}} H(s) \e^{s t}
    \qquad\qquad
    z^{n} \xlongrightarrow{\text{a DT LTI system}} H(z) z^n
\end{equation*}
where $\displaystyle H(s) :=
\int_{-\infty}^{\infty} h(\tau) \e^{-s \tau} \dd\tau$ or
$\displaystyle H(z) := \sum_{k=-\infty}^{\infty}h[k]z^{-k}$
is called the \emph{transfer} or
\emph{system function} of the system.
\begin{proof}
    \begin{equation*}
        h(t) \ast \left(\e^{s t}\right) = 
        \int_{-\infty}^{\infty} h(\tau) \e^{s (\tau - t)} \dd\tau =
        \e^{s t} \int_{-\infty}^{\infty} h(\tau) \e^{-s \tau} \dd\tau
    \end{equation*}
    where $H(s)$ is the eigenvalue.
\end{proof}

Therefore, we have
\begin{equation*}
    x(t)=\sum_i A_i \e^{s_i t} 
    \xlongrightarrow{\text{an LTI system}}
    y(t)=\sum_i H(s_i) A_i \e^{s_i t}
\end{equation*}
In the following sections,
we only focus on the purely imaginary values of $s$, i.e., $s = \ju\omega$.

\section{FS Repr of CT Periodic Signals}

All CT complex exponential
signals that are periodic with period $T$ is given by
\begin{equation*}
    \phi_k(t) := \e^{\ju k \omega_0 t} = \e^{\ju k \frac{2\pi}{T} t}
    \qquad k \in \mathbb{Z}
\end{equation*}
where $\omega_0 = 2\pi/T$.

Many periodic signals can be represented as a linear combination of harmonically
related complex exponentials:
\begin{equation*}
    x(t) = \sum_{k=-\infty}^{\infty} a_k \e^{\ju k \omega_0 t}
\end{equation*}
We choose harmonically related complex exponentials due to their orthogonality:
\begin{theorem box}
    \begin{equation*}
    \left< \e^{\ju m \omega_0 t}, \e^{\ju n \omega_0 t} \right> =
    \int_T \e^{\ju(m-n) \omega_0 t}\dd{t} =
    \mycases{
        T & \qif m=n \\
        0 & \qif m\neq n
        }
    \end{equation*}
    where $m,n\in\mathbb{Z}$ and $T=\dfrac{2\pi}{\omega_0}>0$.
\end{theorem box}

\begin{definition}[The FS of a CT Periodic Signal]
    If a signal $x(t)$ with a period $T$ has a Fourier series representation,
    then
    \begin{gather*}
        x(t) = \sum_{k=-\infty}^{\infty} a_k \e^{\ju k \omega_0 t}
        = \sum_{k=-\infty}^{\infty} a_k \e^{\ju k\frac{2\pi}{T}t}
        \tag{The Synthesis Equation} \\
        a_k = \frac{1}{T} \int_T x(t) \e^{-\ju k\omega_0 t} \dd{t}
        = \frac{1}{T} \int_T x(t) \e^{-\ju k\frac{2\pi}{T}t} \dd{t}
        \tag{The Analysis Equation}
    \end{gather*}
    where the set $\left\{a_k\right\}$ is called
    the \emph{Fourier series coefficients} or
    the \emph{spectral coefficients} of $x(t)$.
\end{definition}

\begin{remark}
    $a_0$ is simply the average value over one period,
    called the \emph{dc} or \emph{constant component} of $x(t)$.
    \begin{equation*}
        a_0 = \frac{1}{T}\int_T x(t) \dd{t}
    \end{equation*}
\end{remark}

The coefficients can be seen as projections of the signal onto the
eigenfunctions of the system:
\begin{theorem box}
\begin{equation*}
    a_k = \frac{\displaystyle\left<x(t), \e^{\ju k \omega_0 t}\right>}
    {\displaystyle\left<\e^{\ju k \omega_0 t}, \e^{\ju k \omega_0 t}\right>} =
    \frac{\displaystyle\int_T x(t) \e^{-\ju k\omega_0 t} \dd{t}}
    {\displaystyle\int_T \e^{\ju k \omega_0 t} \e^{-\ju k\omega_0 t} \dd{t}}
    = \frac{1}{T} \int_T x(t) \e^{-\ju k\omega_0 t} \dd{t}
\end{equation*}
\end{theorem box}

\section{The FS of a DT Periodic Signal}

All DT complex exponential
signals that are periodic with period $N$ is given by
\begin{equation*}
    \phi_k[n] := \e^{\ju k \omega_0 n} = \e^{\ju k \frac{2\pi}{N} n}
    \qquad k \in \mathbb{Z}
\end{equation*}
where $\omega_0 = 2\pi/N$.

However, there are only $N$ distinct signals given by this, because
\begin{theorem box}
    \begin{equation*}
        \phi_{k+N}[n] = \e^{\ju k \frac{2\pi}{N} n + \ju k 2\pi}
        = \e^{\ju k \frac{2\pi}{N} n} = \phi_k[n]
        \qquad \forall k \in \mathbb{Z}
    \end{equation*}
\end{theorem box}

For a single $\phi_k$, the following property holds:
\begin{theorem box}
    \begin{equation*}
        \sum_{n=\langle N\rangle} \phi_k[n] = 
        \sum_{n=\langle N\rangle} \e^{\ju k \frac{2\pi}{N} n} = 
        \mycases{
            N & \qfor k=0,\pm N, \pm 2N, \cdots \\
            0 & \qotherwise
        }
    \end{equation*}
\end{theorem box}
\begin{proof}
    \begin{equation*}
        \sum_{n=\langle N\rangle} \e^{\ju k \frac{2\pi}{N} n}
        = \sum_{n=n_0}^{n_0+N-1} \e^{\ju k \frac{2\pi}{N} n}
        = \mycases{
            \displaystyle
            \sum_{n=\langle N\rangle} 1
            = N & \qfor \e^{\ju k \frac{2\pi}{N}}=1 \\\\
            \displaystyle
            \frac{\displaystyle \e^{\ju k \frac{2\pi}{N} n_0}\left(
                1-\e^{\ju k \frac{2\pi}{N}N}\right)}
            {\displaystyle 1-\e^{\ju k \frac{2\pi}{N}}}
            = 0 & \qotherwise
        }
    \end{equation*}
\end{proof}

\begin{remark}
    ``$\sum\limits_{\scriptscriptstyle n=\langle N\rangle}$''
    means the sum over any interval of $N$ integers.
\end{remark}

\paragraph{Orthogonality}
\begin{equation*}
    \left\langle \phi_k[n], \phi_l[n] \right\rangle
    = \sum_{n=\langle N\rangle}\phi_k[n]\phi_l^\ast[n]
    = \sum_{n=\langle N\rangle}\e^{\ju (k-l)\omega_0 n}
    = \mycases{
        N & \qfor k=l \\
        0 & \qfor k\neq l
    }
\end{equation*}
where $k,l$ are in the same period, i.e., $\abs{k-l}<N$.

\begin{definition}[The FS of a DT Periodic Signal]
    If a signal $x[n]$ with a period $N$ has a Fourier series representation,
    then
    \begin{gather*}
        x[n] = \sum_{k=\langle N\rangle} a_k \e^{\ju k \omega_0 n}
        = \sum_{k=\langle N\rangle} a_k \e^{\ju k\frac{2\pi}{T}n}
        \tag{The Synthesis Equation} \\
        a_k = \frac{1}{N} \sum_{n=\langle N\rangle} x[n] \e^{-\ju k\omega_0 n}
        = \frac{1}{N} \sum_{n=\langle N\rangle} x[n] \e^{-\ju k\frac{2\pi}{N}n}
        \tag{The Analysis Equation}
    \end{gather*}
    where the set $\left\{a_k\right\}$ is called
    the \emph{Fourier series coefficients} or
    the \emph{spectral coefficients} of $x[n]$.
\end{definition}
The coefficients are also periodic with period $N$, i.e.,
\begin{theorem box}
    \begin{equation*}
        a_k = a_{k+N} \qquad \forall k \in \mathbb{Z}
    \end{equation*}
\end{theorem box}

\section{Properties of the Fourier Series}

\subsection{}

Suppose we have $x(t)$ and $y(t)$, or $x[n]$ and $y[n]$,
of the \textbf{same period} $T$ or $N$, and
\begin{gather*}
    \mycases[0.7]{
        x(t) \xrightleftharpoons[\mathcal{FS}^{-1}]{\mathcal{FS}} a_k \\
        y(t) \xrightleftharpoons[\mathcal{FS}^{-1}]{\mathcal{FS}} b_k
    } \qor
    \mycases[0.7]{
        x[n] \xrightleftharpoons[\mathcal{FS}^{-1}]{\mathcal{FS}} a_k \\
        y[n] \xrightleftharpoons[\mathcal{FS}^{-1}]{\mathcal{FS}} b_k
    }
\end{gather*}
Then we have the following results:
\begin{align*}
\shortintertext{Linearity}
    A\cdot x(t) + B\cdot y(t)
    \quad&\xrightleftharpoons[\mathcal{FS}^{-1}]{\mathcal{FS}}\quad
    A\cdot a_k + B\cdot b_k
    &
    A\cdot x[n] + B\cdot y[n]
    \quad&\xrightleftharpoons[\mathcal{FS}^{-1}]{\mathcal{FS}}\quad
    A\cdot a_k + B\cdot b_k
\shortintertext{Time Shifting}
    x(t-t_0)
    \quad&\xrightleftharpoons[\mathcal{FS}^{-1}]{\mathcal{FS}}\quad
    \e^{-\ju k\omega_0 t_0} a_k
    &
    x[n-n_0]
    \quad&\xrightleftharpoons[\mathcal{FS}^{-1}]{\mathcal{FS}}\quad
    \e^{-\ju k\omega_0 n_0} a_k
\shortintertext{Frequency Shifting}
    \e^{\ju M \omega_0 t} x(t)
    \quad&\xrightleftharpoons[\mathcal{FS}^{-1}]{\mathcal{FS}}\quad
    a_{k-M}
    &
    \e^{\ju M \omega_0 n} x[n]
    \quad&\xrightleftharpoons[\mathcal{FS}^{-1}]{\mathcal{FS}}\quad
    a_{k-M}
\shortintertext{Time Reversal}
    x(-t)
    \quad&\xrightleftharpoons[\mathcal{FS}^{-1}]{\mathcal{FS}}\quad
    a_{-k}
    &
    x[-n]
    \quad&\xrightleftharpoons[\mathcal{FS}^{-1}]{\mathcal{FS}}\quad
    a_{-k}
\shortintertext{Time Scaling}
    x(\alpha t)
    \quad&\xrightleftharpoons[\mathcal{FS}^{-1}]{\mathcal{FS}}\quad
    a_k \begin{pmatrix}
        \text{with fundamental} \\
        \text{frequency } \alpha\omega_0
    \end{pmatrix}
    &
    \mycases{
        x[\frac{n}{m}] & \text{for } \frac{n}{m}\in\mathbb{Z} \\
        0 & \text{otherwise}
    }
    \quad&\xrightleftharpoons[\mathcal{FS}^{-1}]{\mathcal{FS}}\quad
    \frac{1}{m}a_k \begin{pmatrix}
        \text{with period} \\
        \text{of } m N
    \end{pmatrix}
\shortintertext{Conjugation}
    x^\ast(t)
    \quad&\xrightleftharpoons[\mathcal{FS}^{-1}]{\mathcal{FS}}\quad
    a_k^\ast
    &
    x^\ast[n]
    \quad&\xrightleftharpoons[\mathcal{FS}^{-1}]{\mathcal{FS}}\quad
    a_k^\ast
\shortintertext{Periodic Convolution}
    \int_T x(\tau) y(t-\tau)\dd{\tau}
    \quad&\xrightleftharpoons[\mathcal{FS}^{-1}]{\mathcal{FS}}\quad
    T a_k b_k
    &
    \sum_{r=\langle N\rangle} x[r] y[n-r]
    \quad&\xrightleftharpoons[\mathcal{FS}^{-1}]{\mathcal{FS}}\quad
    N a_k b_k
\shortintertext{Multiplication}
    \begin{pmatrix}
        \text{of same} \\
        \text{period}
    \end{pmatrix}
    x(t) y(t)
    \quad&\xrightleftharpoons[\mathcal{FS}^{-1}]{\mathcal{FS}}\quad
    a_k \ast b_k
    &
    x[n] y[n]
    \quad&\xrightleftharpoons[\mathcal{FS}^{-1}]{\mathcal{FS}}\quad
    \sum_{r=\langle N\rangle} a_r b_{k-r}
\shortintertext{Differentiation \& Difference}
    \dv{x(t)}{t}
    \quad&\xrightleftharpoons[\mathcal{FS}^{-1}]{\mathcal{FS}}\quad
    \left(\ju k\omega_0\right) a_k
    &
    x[n] - x[n-1]
    \quad&\xrightleftharpoons[\mathcal{FS}^{-1}]{\mathcal{FS}}\quad
    \left(1-\e^{-\ju \omega_0}\right) a_k
\shortintertext{Integration \& Running Sum}
    \int_{-\infty}^{t}x(\tau)\dd{\tau}
    \quad&\xrightleftharpoons[\mathcal{FS}^{-1}]{\mathcal{FS}}\quad
    \frac{1}{\ju k\omega_0} a_k
    &
    \sum_{k=-\infty}^{n} x[k]
    \quad&\xrightleftharpoons[\mathcal{FS}^{-1}]{\mathcal{FS}}\quad
    \frac{1}{1-\e^{-\ju k\omega_0}} a_k
\shortintertext{Parseval's Relation for Periodic Signals}
    \frac{1}{T} \int_T \abs{x(t)}^2 \dd{t}
    \quad&=\quad
    \sum_{k=-\infty}^{\infty} \abs{a_k}^2
    &
    \frac{1}{N} \sum_{n=\langle N\rangle} \abs{x[n]}^2
    \quad&=\quad
    \sum_{k=\langle N\rangle} \abs{a_k}^2
\end{align*}

\subsection{Symmetry}

Both CTFS and DTFS have exactly the same symmetry properties as the CTFT,
including conjugate symmetry, parity, and even-odd decomposition.

\section{Transfer Functions}

For a given system, $H$ is called the \textbf{transfer function},
\textbf{system function}, or \textbf{frequency response}.

\subsection{CT Periodic Signals}
\begin{definition}[Transfer Function of a CT System]
    Suppose $h(t)$ is the \emph{impulse response} of the system,
    then the transfer function $H$ is given by
    \begin{equation*}
        H(\ju\omega) := \int_{-\infty}^{\infty} h(t) \e^{-\ju\omega t} \dd{t}
    \end{equation*}
\end{definition}

\begin{theorem box}
    Suppose $x(t)$ is a periodic signal and
    $a_k$ is the FS representation of $x(t)$.
    \begin{gather*}
        y(t) = x(t) \ast h(t)
        \quad\xlongrightarrow{\mathcal{FS}}\quad
        a_k\cdot H(\ju k\omega_0)
    \end{gather*}
\end{theorem box}
\begin{proof}
    \begin{gather*}
        x(t) = \sum_{k=-\infty}^{\infty} a_k \e^{\ju k \omega_0 t}
        \quad\xlongrightarrow{\mathcal{FS}}\quad
        y(t) = \sum_{k=-\infty}^{\infty} a_k H(\ju k\omega_0)\e^{\ju k \omega_0 t}
    \end{gather*}
\end{proof}

\begin{definition}[Transfer Function of a DT System]
    Suppose $h[n]$ is the \emph{impulse response} of the system,
    then the transfer function $H$ is given by
    \begin{equation*}
        H(\e^{\ju\omega}) := \sum_{n=-\infty}^{\infty} h[n] \e^{-\ju\omega n}
    \end{equation*}
\end{definition}

\begin{theorem box}
    Suppose $x[n]$ is a periodic signal and
    $a_k$ is the FS representation of $x[n]$.
    ($a_k$ is periodic.)
    \begin{gather*}
        y[n] = x[n] \ast h[n]
        \quad\xlongrightarrow{\mathcal{FS}}\quad
        a_k\cdot H(\e^{\ju k\omega_0})
    \end{gather*}
\end{theorem box}

\end{document}